\section{Εισαγωγή στο \ENG{E-Learning} και την Ανάγκη Προτυποποίησης}
\label{sec:elearning_intro}

Η ηλεκτρονική μάθηση (\textit{\ENG{E-Learning}}) αποτελεί σημαντική καινοτομία στον τομέα της εκπαίδευσης και της 
επαγγελματικής κατάρτισης κατά τη διάρκεια των τελευταίων δεκαετιών. Ο όρος \textit{\ENG{E-Learning}} αναφέρεται
στη χρήση ηλεκτρονικών μέσων, κυρίως διαδικτυακών τεχνολογιών και ψηφιακών πλατφορμών, για την παροχή εκπαιδευτικού περιεχομένου,
την υποστήριξη της μαθησιακής διαδικασίας και την παρακολούθηση της προόδου των εκπαιδευομένων. Το \ENG{E-Learning} φαίνεται
να πρωτοεμφανίστηκε στη δεκαετία του 1960, με τα τα πρώτα συστήματα υπολογιστικής εκπαίδευσης (\textit{\ENG{Computer-Based} \ENG{Training},
\ENG{CBT}}), τα οποία αξιοποιούσαν μαγνητικές ταινίες, δισκέτες και αργότερα οπτικούς δίσκους (\ENG{CD-ROMs}) για την αποθήκευση και διανομή
εκπαιδευτικού υλικού.

Κατά τη δεκαετία του 1980 και στις αρχές της δεκαετίας του 1990, η ανάπτυξη εκπαιδευτικού λογισμικού και περιεχομένου χαρακτηριζόταν από έντονη 
κατακερματισμότητα και έλλειψη διαλειτουργικότητας. Κάθε προμηθευτής λογισμικού (\textit{\ENG{vendor}}) ανέπτυσσε ιδιόκτητες (\textit{\ENG{proprietary}})
λύσεις, οι οποίες ήταν ασυμβίβαστες μεταξύ τους. Το εκπαιδευτικό περιεχόμενο που δημιουργούνταν για μία συγκεκριμένη πλατφόρμα δεν μπορούσε να εκτελεστεί
σε άλλη, επιβάλλοντας στους οργανισμούς να επενδύουν σημαντικούς πόρους για την ανάπτυξη ή προσαρμογή του περιεχομένου κάθε φορά που άλλαζαν σύστημα
διαχείρισης μάθησης (\textit{\ENG{Learning} \ENG{Management} \ENG{System}, \ENG{LMS}}). Αυτή η κατάσταση οδηγούσε σε υψηλό οικονομικό κόστος,
περιορισμένη επαναχρησιμοποίηση εκπαιδευτικών πόρων και τεχνολογική εξάρτηση από συγκεκριμένους προμηθευτές (\textit{\ENG{vendor} \ENG{lock-in}}).

Η ανάγκη για διαλειτουργικότητα (\textit{\ENG{interoperability}}), επαναχρησιμοποιησιμότητα (\textit{\ENG{reusability}}),
προσβασιμότητα (\textit{\ENG{accessibility}}) και μεταφερσιμότητα (\textit{\ENG{portability}}) του εκπαιδευτικού περιεχομένου έγινε εμφανής
στα μέσα της δεκαετίας του 1990. Οι οργανισμοί, ιδίως στον τομέα της άμυνας, της αεροπορίας και των μεγάλων επιχειρήσεων, αντιμετώπιζαν
σημαντικές δαπάνες για την ανάπτυξη και συντήρηση εκπαιδευτικών συστημάτων. Η έλλειψη κοινών προδιαγραφών οδηγούσε σε επανάληψη ενεργιών, 
αδυναμία ανταλλαγής περιεχομένου μεταξύ οργανισμών και δυσκολία στην αξιολόγηση της αποτελεσματικότητας των εκπαιδευτικών προγραμμάτων.

Σύμφωνα με τη βιβλιογραφία, το οικονομικό κόστος της μη-προτυποποίησης εκτιμάται ότι μπορούσε να φτάσει το 30-50\% του συνολικού προϋπολογισμού
ανάπτυξης εκπαιδευτικού περιεχομένου, καθώς κάθε μετάβαση σε νέα πλατφόρμα απαιτούσε επανασχεδιασμό και επαναδημιουργία του περιεχομένου.
Επιπλέον, η αδυναμία συλλογής και σύγκρισης δεδομένων προόδου μεταξύ διαφορετικών συστημάτων περιόριζε τις δυνατότητες αξιολόγησης και βελτίωσης
των εκπαιδευτικών προγραμμάτων.

Η προτυποποίηση αναδείχθηκε ως λύση στα παραπάνω προβλήματα. Τα πρότυπα (\textit{\ENG{standards}}) ορίζουν κοινές προδιαγραφές για τη δομή,
την αποθήκευση, την επικοινωνία και την παρακολούθηση του εκπαιδευτικού περιεχομένου, επιτρέποντας σε διαφορετικά συστήματα να συνεργάζονται
και να ανταλλάσσουν πληροφορίες.

Στο πλαίσιο αυτό, τα \ENG{E-Learning} \ENG{standards} βασίζονται σε ορισμένες θεμελιώδεις αρχές.
Πρώτον, τη διαλειτουργικότητα (\ENG{interoperability}), δηλαδή την ικανότητα ενός εκπαιδευτικού πακέτου να εκτελείται σε διαφορετικά
\ENG{LMS} χωρίς τροποποιήσεις. Δεύτερον, την επαναχρησιμοποιησιμότητα (\ENG{reusability}), που επιτρέπει 
το ίδιο εκπαιδευτικό περιεχόμενο να αξιοποιείται σε διαφορετικά εκπαιδευτικά πλαίσια και προγράμματα. 
Τρίτον, τη μεταφερσιμότητα (\ENG{portability}), η οποία αναφέρεται στην ευκολία μεταφοράς περιεχομένου από ένα σύστημα σε άλλο.
Παράλληλα, ιδιαίτερη σημασία έχει η προσβασιμότητα (\ENG{accessibility}), ώστε το εκπαιδευτικό υλικό να είναι διαθέσιμο και 
σε άτομα με ειδικές ανάγκες. Τέλος, σημαντικό ρόλο διαδραματίζει η ανευρεσιμότητα (\ENG{discoverability}),
μέσω της χρήσης μεταδεδομένων που επιτρέπουν την αποτελεσματική αναζήτηση και ταξινόμηση εκπαιδευτικών πόρων.

Η προτυποποίηση στον τομέα του \ENG{E-Learning} αποτέλεσε μια πολυεπίπεδη και πολυετή προσπάθεια, στην οποία συμμετείχαν κυβερνητικοί οργανισμοί (ιδίως το Υπουργείο Άμυνας των ΗΠΑ), βιομηχανικές ενώσεις (όπως η \ENG{Aviation} \ENG{Industry} \ENG{Computer-Based} \ENG{Training} \ENG{Committee}), ακαδημαϊκά ιδρύματα και εταιρείες τεχνολογίας. Στις επόμενες ενότητες, εξετάζονται αναλυτικά τα κύρια στάδια αυτής της εξέλιξης, από την εποχή πριν την προτυποποίηση μέχρι τα σύγχρονα πρότυπα που κυριαρχούν σήμερα.

\begin{figure}[H]
  \centering
  % TikZ-based horizontal timeline
\begin{chapterfigurebox}
\begin{tikzpicture}[
  scale=0.85,
  every node/.style={font=\small},
  timeline/.style={draw, thick, -latex},
  era/.style={rectangle, rounded corners, draw, fill=white, minimum width=2cm, minimum height=0.8cm, align=center, drop shadow},
  milestone/.style={circle, draw, fill=white, inner sep=2pt, font=\tiny\bfseries}
]

% Main timeline axis
\draw[timeline, line width=2pt] (0,0) -- (16,0) node[right, font=\small\bfseries] {2026};

% Timeline markers (years)
\foreach \x/\year in {0/1980, 2/1988, 4/1995, 6/2000, 8/2004, 10/2010, 12/2016, 14/2020, 16/2026} {
  \draw[thick] (\x,0.1) -- (\x,-0.1) node[below, font=\scriptsize] {\year};
}

% ERA 1: Pre-Standardization (1980-1993)
\node[era, fill=problemred!20, text width=2.5cm] at (1, 2) (pre) {
  \textbf{Προ-\ENG{Standard}}\\
  \tiny \ENG{CBT}\\
  \tiny Ιδιόκτητες λύσεις
};
\draw[->, thick, problemred] (pre) -- (1, 0.2);

% ERA 2: AICC Foundation (1988)
\node[era, fill=aiccblue!30] at (2, -2) (aicc) {
  \textbf{\ENG{AICC}}\\
  \tiny 1988
};
\draw[->, thick, aiccblue] (aicc) -- (2, -0.2);

% ERA 3: Metadata Standards (1995-2002)
\node[era, fill=metadatagray!30, text width=2cm] at (4, 2.2) (metadata) {
  \textbf{\ENG{Metadata}}\\
  \tiny \ENG{Dublin} \ENG{Core}\\
  \tiny \ENG{IEEE} \ENG{LOM}
};
\draw[->, thick, metadatagray] (metadata) -- (4, 0.2);

% ERA 4: IMS Global (1997)
\node[era, fill=imslightblue!40, text width=1.8cm] at (5, -2.2) (ims) {
  \textbf{\ENG{IMS} \ENG{Global}}\\
  \tiny 1997
};
\draw[->, thick, imslightblue] (ims) -- (5, -0.2);

% ERA 5: SCORM Birth (2000-2004)
\node[era, fill=scormgreen!40, text width=2cm] at (7, 2.5) (scorm) {
  \textbf{\ENG{SCORM}}\\
  \tiny 1.0 → 2004
};
\draw[->, thick, scormgreen] (scorm) -- (7, 0.2);

% ERA 6: SCORM Maturity (2004-2009)
\node[era, fill=scormgreen!60, text width=2cm] at (9, -2.5) (scorm2) {
  \textbf{\ENG{SCORM} 2004}\\
  \tiny 4\ENG{th} \ENG{Edition}
};
\draw[->, thick, scormgreen!80] (scorm2) -- (9, -0.2);

% ERA 7: xAPI Revolution (2011-2016)
\node[era, fill=xapiorange!50, text width=2cm] at (11, 2.8) (xapi) {
  \textbf{\ENG{xAPI}}\\
  \tiny \ENG{Tin} \ENG{Can} \ENG{API}\\
  \tiny 2011
};
\draw[->, thick, xapiorange] (xapi) -- (11, 0.2);

% ERA 8: cmi5 (2016)
\node[era, fill=accentpurple!40] at (12, -2.8) (cmi5) {
  \textbf{\ENG{cmi5}}\\
  \tiny 2016
};
\draw[->, thick, accentpurple] (cmi5) -- (12, -0.2);

% ERA 9: Contemporary (2020+)
\node[era, fill=blue!20, text width=2.5cm] at (15, 3) (contemporary) {
  \textbf{\ENG{AI/Analytics} \ENG{Era}}\\
  \tiny \ENG{LRS}, \ENG{xAPI} 2.0\\
  \tiny \ENG{Generative} \ENG{AI}
};
\draw[->, thick, blue!70] (contemporary) -- (15, 0.2);

% Key milestones (dots)
\node[milestone, fill=aiccblue] at (2,0) {};
\node[milestone, fill=scormgreen] at (6,0) {};
\node[milestone, fill=xapiorange] at (11,0) {};
\node[milestone, fill=accentpurple] at (12,0) {};

% Legend
\node[below, font=\tiny, align=left] at (8, -3.5) {
  \textcolor{problemred}{\textbullet} Προ-\ENG{Standard} (1980-1993) \quad
  \textcolor{aiccblue}{\textbullet} \ENG{AICC/IMS} (1988-2000) \quad
  \textcolor{scormgreen}{\textbullet} \ENG{SCORM} \ENG{Era} (2000-2010) \quad
  \textcolor{xapiorange}{\textbullet} \ENG{xAPI/cmi5} (2011-σήμερα)
};

\end{tikzpicture}
\end{chapterfigurebox}

  \caption{Χρονολογική Εξέλιξη των \ENG{E-Learning} \ENG{Standards} (1980-2026)}
  \label{fig:timeline_evolution}
\end{figure}
